\textit{Night-observation main variant.}
You have 30 minutes in total. The telescope and finder tube are already aligned.
Do not change the alignment. Write answers on the separate answer sheet.
Answers and calculations on the question paper are not evaluated.

\begin{description}
\item[(a)] What is the age of the Moon today? The age is the number of days
since the last new Moon. \hfill\mbox{[5 points]}
\item[(b)] What is the angular distance between the Moon and Saturn?
\hfill\mbox{[5 points]}
\item[(c)] After how many days will the Moon be close to Saturn?
\hfill\mbox{[10 points]}
\item[(d)] What is the altitude of the Sun at the time of observation? After
answering, give the examiner your answer sheet to record the time.
\hfill\mbox{[10 points]}
\item[(e)] Point the telescope at the Moon and centre it in the field of view.
Ask the examiner to confirm successful targeting and record it on your answer
sheet. If you cannot target the Moon, the examiner will do so and you lose
the points for this part. \hfill\mbox{[5 points]}
\item[(f)] The image shows the Moon with Mare Tranquilitatis marked with an
\(\times\). Ask the examiner for a stopwatch and use it to measure the angular
diameter of the mare. If you lose the Moon and cannot retarget it, give the
examiner your answer sheet and ask for help. You lose any points earned in (e).
\hfill\mbox{[10 points]}
\end{description}
\ioaafig{2023_OBS-NIGHT_Lunar-f1.pdf}
